
\section{Introduction}

% Classical TN contractions
Tensor network contractions have been identified as central inference mechanism of neuro-symbolic AI.
However, their exact execution is $\mathrm{NP}$-hard, and one either has to restrict itself to feasible situations, or apply approximations.

% Quantum
We in this chapter investigate, how specific contractions (especially forward inference on \ComputationActivationNetworks{}) can be performed more efficiently on a Quantum Computer.

Tensor Network contractions can be performed by quantum computing as follows:
\begin{itemize}
    \item \textbf{Rejection Sampling based:}
    Can augment any tensor network by an Bayesian network, which is efficiently inferred by forward sampling.
    We prepare quantum states, which measurement statistics can be utilized to prepare samples from interesting distributions such as \ComputationActivationNetworks{}.
    The square root quantum advantage is obtained by amplitude amplification.
    \item \textbf{Phase Estimation based:}
    Measure the eigenvalues of a Grover operator to a formula and a q-sample.
    \item \textbf{Fourier Transform based:}
    Limited contractions (given promises on their possible values) can be formulated as Hidden Subgroup Problems and solved by Fourier transforms.
%    Here we investigate how we can exploit these as contraction provider for \tnreason{}.
    We are inspired by Deutsch-Jozsa algorithm, which we generalize here to a general moment computation problem.
\end{itemize}

%
Note that the rejection sampling scheme and the phase estimation are ignorant under the phase of the prepared state, and only the Fourier transform based scheme depend on the phase.

% NP Hardness of contraction:
Generic tensor network contraction is $\mathrm{NP}$-hard, reflecting the hardness of logical and probabilistic inference.
Quantum Computers can prepare states, which measurement distribution are specific network contractions, which can be hard to contract classically.
We study schemes to retrieve information about these states, for example by sampling schemes, and use the retrieved information in inference.

\subsection{Circuit Preparation schemes}

Circuit preparing schemes based on approximation:
\begin{itemize}
    \item Q-Alchemy \cite{araujo_low-rank_2023}, Q-Tucker \cite{blank_tucker_2026}
    \item Tensor-Network Optimization based (alternating schemes) \cite{rudolph_synergistic_2023,rudolph_decomposition_2023}
    \item Preparing a quantum circuit given an MPS requires logarithmic depth in the tensor order \cite{malz_preparation_2024}
\end{itemize}

Exact circuit preparing schemes for distributions:
\begin{itemize}
    \item Grover-Rudolph \cite{grover_creating_2002}, but no quantum advantage \cite{herbert_no_2021}
    \item Uniform controlled rotations \cite{mottonen_transformation_2005}
    \item Quantum Shannon decomposition \cite{shende_synthesis_2006}
\end{itemize}

Circuit simulation:
Since \qcreason{} can prepare quantum circuits to arbitrary tensor networks, it can also be used to simulate quantum circuits (with an overhead!).
\begin{itemize}
    \item \cite{sander_large-scale_2025,sander_quantum_2025}
\end{itemize}

\subsection{Related Works}

Quantum algorithms for linear algebra:
\begin{itemize}
    \item Contraction by SWAP and Hadamard test (see Appendix)
    \item Deutsch-Jozsa Algorithm
    \item HHL algorithm to solve linear equations \cite{harrow_quantum_2009}
\end{itemize}